\documentclass[11pt]{article}
\usepackage[margin=1in]{geometry}
\usepackage{booktabs}
\title{DME: SNc vs.\ VTA, within Lesioned Cells --- Interpretation}
\author{GSE233866, 6-OHDA lesion cohort}
\date{}

\begin{document}
\maketitle

This is the most disease-\emph{and}-region-specific comparison in the
project: not baseline regional structure, and not disease-vs-healthy within
one region, but which genes and modules are differentially active in SNc
versus VTA specifically while neurodegeneration is in progress --- the
signal most directly relevant to a region-selective treatment target.
SNc-lesioned (2{,}291 cells) and VTA-lesioned (1{,}720 cells) are pooled into
one 4{,}011-cell network before testing.

CACNA1D's module here is \textbf{blue} ($\mathrm{kME}=0.32$), and it shows
one of the strongest signals in the entire table: $\mathrm{avg\_log2FC} =
+13.34$, $p_{\mathrm{adj}} = 4.5\times10^{-123}$, higher in SNc. Eleven of
twelve tested modules differ significantly between the two regions during
lesioning, which is expected on its own --- SNc and VTA are fundamentally
different neuron populations, not disease-state variants of the same one ---
but CACNA1D's own module being among the most extreme values in that set is
notable specifically because it is CACNA1D's module, not an arbitrary one.

Taken together with the lesioned-vs-intact result in
\texttt{lesioned\_vs\_intact\_combined\_dme\_snc/}, this is the most direct
evidence this project has produced for something CACNA1D-network-specific
happening preferentially in the vulnerable region \emph{during} the
neurodegenerative process itself, rather than as a static baseline
difference. As with the other DME table, treat log2FC magnitude as
directional evidence only; module eigengenes can be negative or near-zero,
which destabilizes fold-change arithmetic even when the underlying
$p$-value is sound.

\end{document}
